\chapter{Phonology: The Glottalic Architecture}

\section{The Ancestral Ejective Series}
The foundational phonological invariant of PED is the existence of a series of three ejective stops: \textit{*P', *T', *K'}. Mainstream Indo-European studies have traditionally posited a series of 'voiced aspirates' (\textit{*bh, *dh, *gh}), but as Hopper (1973) and Gamkrelidze (1973) argued, this creates a typologically unstable system. The PED framework recognizes that these were originally harsh, glottalized sounds that softened differently across the Eurasian landmass.

\section{The Dental Matrix (*T')}
The dental ejective \textit{*T'} is the most stable and heavily attested phoneme in our reconstruction. Its divergence follows a strictly deterministic path.

\subsection{The Western Node (Vedic Sanskrit)}
In the Western branch, the PED ejective lost its glottalic stricture, surfacing as the plain voiced stop \textit{*d} in Vedic. This is preserved in the core root for 'wood/tree' (\textit{*T'ar-} > Vedic \textit{dāru}). 

\subsection{The Southern Node (Old Tamil)}
In the South, the ejective stricture was lost, but the voiceless quality was retained, surfacing as \textit{*t}. In Old Tamil, this often geminates (\textit{*tt}) to mark transitivity—a structural memory of the original glottal burst. 

\subsection{The Eastern Node (Old Tibetan)}
In the East, the glottalic energy was converted into aspiration, yielding \textit{*th}. This is visible in the Old Tibetan \textit{thig} (mark) and \textit{shing} (wood).

\begin{table}[h]
\centering
\begin{tabular}{@{}lccccl@{}} \toprule
\textbf{PED Root} & \textbf{Vedic} & \textbf{Tamil} & \textbf{Tibetan} & \textbf{Semantic Core} \\ \midrule
*T'ar- & dāru & tāram & shing & Wood/Extension \\
*T'ik- & diś & tikku & thig & Point/Indicate \\
*T'uH- & udan & nīr & chu & Water/Liquid \\
*T'ak- & dagh & takk- & thak & Reach/Enclosure \\ \bottomrule
\end{tabular}
\caption{The Dental Divergence Matrix.}
\end{table}

\section{The Grand Consonantal Matrix}
The following table summarizes the systematic divergence of the PED phonological inventory across the three primary geographic nodes.

\begin{table}[h]
\centering
\begin{tabular}{@{}lccc@{}} \toprule
\textbf{PED Segment} & \textbf{Vedic Reflex} & \textbf{Tamil Reflex} & \textbf{Tibetan Reflex} \\ \midrule
\textbf{*P' (Labial Stop)} & b & p & ph \\
\textbf{*T' (Dental Stop)} & d & t & th \\
\textbf{*K' (Velar Stop)} & g & k & kh \\ \midrule
\textbf{*S (Fricative)} & s & c & s \\
\textbf{*H (Laryngeal)} & h/0 & v/0 & x \\ \midrule
\textbf{*M (Bilabial Nasal)} & m & m & m \\
\textbf{*N (Alveolar Nasal)} & n & n & n \\ \bottomrule
\end{tabular}
\caption{The Full PED Consonantal Matrix.}
\end{table}

\section{The Laryngeal Force (*H)}
The ancestral pharyngeal fricative \textit{*H} is responsible for the 'vowel coloring' that defines PIE and the 'tonogenesis' that defines PST.

\subsection{Laryngeal Scarring Analysis}
\begin{table}[h]
\centering
\begin{tabular}{llll} \toprule
\textbf{Feature} & \textbf{Vedic Reflex} & \textbf{Tamil Reflex} & \textbf{Tibetan Reflex} \\ \midrule
Primary "Scar" & Vowel Coloring & Length/Gemination & Tonogenesis \\
Phonetic Relic & Visarga (\d{h}) & \={A}ytam (ஃ) & Glottal Coda (') \\
Vowel Change & *eH > \={a} & *aH > \={a} & (Tone creation) \\ \bottomrule
\end{tabular}
\caption{The Positional Impact of *H.}
\end{table}

\section{Phonetic Stochasticity and Analogical Leveling}
As noted in Chapter 2, our model incorporates a 15\% exception rate. These are not failures of the law, but 'analogical leveling'—the process by which speakers regularize common words. For instance, the PED root for 'stone' (*K'aL-) surfaces as \textit{kal} in Tamil but \textit{rdo} in Tibetan, demonstrating a lexical replacement that overrides the phonological law. 
\begin{center}
\textit{*K'aL- (Stone) $\rightarrow$ Tamil kal, Tibetan rdo, Vedic aśman.}
\end{center}
Such replacements are expected in any 40,000-year-old system.
